% topic: algorithms
\question หยวนมีไข่ที่เหมือนกันทุกประการจำนวน \textbf{2 ฟอง} และต้องการทดสอบ\textbf{ความทนทาน}ของไข่ โดยมีอาคารสูง \textbf{15 ชั้น} (ชั้นที่ 1 ถึง 15) หยวนสามารถนำไข่ไปปล่อยจากชั้นต่าง ๆ เพื่อสังเกตว่าไข่แตกหรือไม่
\newline
\textbf{นิยาม:} ไข่ทุกฟองมีค่าความทนทาน $T$ เท่ากัน โดย $T$ เป็นจำนวนเต็มตั้งแต่ $1$ ถึง $15$
\begin{itemize}
  \item ปล่อยจากชั้น $\le T$ แล้ว ไข่\textbf{ไม่แตก} (ใช้ซ้ำได้)
  \item ปล่อยจากชั้น $> T$ แล้ว ไข่\textbf{แตก} (ใช้ต่อไม่ได้)
\end{itemize}

\textbf{กติกา:}
\begin{itemize}
  \item มีไข่เพียง 2 ฟอง ถ้าแตกทั้ง 2 ฟองโดยยังระบุ $T$ ไม่ได้ = ล้มเหลว
  \item หยุดการทดสอบเมื่อ\textbf{ระบุค่า $T$ ได้แน่นอน}
\end{itemize}

\textbf{ตัวอย่างกลยุทธ์:} ปล่อยไข่ฟองแรกที่ชั้น 8
\begin{itemize}
  \item ถ้าไข่แตกที่ชั้น 8: $T \le 7$ ใช้ไข่ฟองที่ 2 ทดสอบชั้น 1, 2, 3, …, 7 ตามลำดับ ในกรณีเลวร้ายที่สุด (ไข่ฟองที่ 2 แตกที่ชั้น 7) จะทดสอบทั้งหมด $1 + 7 = 8$ ครั้ง
  \item ถ้าไข่ไม่แตกที่ชั้น 8: $T \ge 8$ ยังมีไข่ 2 ฟอง และเหลือชั้น 9--15 ให้ทดสอบ
\end{itemize}

\textbf{คำถาม:} หยวนต้องการออกแบบกลยุทธ์ที่\textbf{รับประกัน}ว่าสามารถระบุค่า $T$ ได้ในจำนวนครั้งที่น้อยที่สุดในกรณีเลวร้ายที่สุด (Worst Case) จงหาว่าในกรณีที่เลวร้ายที่สุด หยวนจะต้องปล่อยไข่\textbf{อย่างมากที่สุดกี่ครั้ง} (ตอบเป็นตัวเลขจำนวนเต็ม)
\newline\newline
ตอบ ในกระดาษคำตอบ
% answer: 5
% solution:
%   หลักการ: ใช้ไข่ฟองแรกกระโดดข้ามช่วงที่กว้าง แล้วค่อย ๆ ลดช่วงลง
%   เพื่อให้ไข่ฟองที่สองไม่ต้องทดสอบทีละชั้นมากเกินไป
%
%   กลยุทธ์ที่ดีที่สุด: ปล่อยไข่ฟองแรกที่ชั้น 5, 9, 12, 14, 15 (ลดช่วงทีละ 1)
%   ช่วงห่าง: 5, 4, 3, 2, 1 — ครอบคลุม 15 ชั้นพอดี
%
%   ตรวจสอบกรณีเลวร้าย:
%     T=4:  ฟอง1แตกที่ชั้น 5 → ฟอง2ทดสอบชั้น 1-4 (4 ครั้ง)            รวม 1+4=5 ✓
%     T=8:  ฟอง1รอดชั้น 5, แตกที่ชั้น 9 → ฟอง2ทดสอบชั้น 6-8 (3 ครั้ง)   รวม 2+3=5 ✓
%     T=11: ฟอง1รอดชั้น 5,9, แตกที่ชั้น 12 → ฟอง2ทดสอบชั้น 10-11 (2 ครั้ง) รวม 3+2=5 ✓
%     T=13: ฟอง1รอดชั้น 5,9,12, แตกที่ชั้น 14 → ฟอง2ทดสอบชั้น 13 (1 ครั้ง) รวม 4+1=5 ✓
%     T=14: ฟอง1รอดชั้น 5,9,12,14, แตกที่ชั้น 15 → ทราบว่า T=14 ทันที   รวม 5 ครั้ง ✓
%     T=15: ฟอง1รอดทุกชั้น → ทราบว่า T=15                              รวม 5 ครั้ง ✓
%
%   สูตรทั่วไป: ต้องการ k ที่น้อยที่สุดที่ทำให้ k(k+1)/2 ≥ 15
%     k=4: 4×5/2 = 10 < 15  (ไม่พอ)
%     k=5: 5×6/2 = 15 ≥ 15  ✓
%   ดังนั้น k=5 คือจำนวนครั้งในกรณีเลวร้ายที่สุด
%
%   เปรียบเทียบกับกลยุทธ์เริ่มที่ชั้น 8: Worst Case = 8 ครั้ง
%   กลยุทธ์ที่ดีที่สุด: Worst Case = 5 ครั้ง
